\documentclass[11pt]{article}

\usepackage{amsmath,amssymb}
\usepackage{xurl}
\usepackage[colorlinks=true,linkcolor=blue,citecolor=blue,urlcolor=blue]{hyperref}

% Shared commands for notes in docs/paper-gaps/.

\DeclareUnicodeCharacter{2080}{\ifmmode{}_{0}\else\textsubscript{0}\fi}
\DeclareUnicodeCharacter{2081}{\ifmmode{}_{1}\else\textsubscript{1}\fi}
\DeclareUnicodeCharacter{2082}{\ifmmode{}_{2}\else\textsubscript{2}\fi}
\DeclareUnicodeCharacter{2099}{\ifmmode{}_{n}\else\textsubscript{n}\fi}

\newcommand{\Lean}{\textsc{Lean}}
\ExplSyntaxOn
\NewDocumentCommand{\leanid}{m}{
  \ifmmode
    \textsf{\detokenize{#1}}
  \else
    \group_begin:
    \urlstyle{sf}
    \tl_set:Nn \l_tmpa_tl {#1}
    \tl_replace_all:Nnn \l_tmpa_tl {\_} {_}
    \exp_args:NV \path \l_tmpa_tl
    \group_end:
  \fi
}
\ExplSyntaxOff
\providecommand{\tr}{\operatorname{tr}}
\newcommand{\tnleanIssueBase}{https://github.com/LionSR/TNLean/issues/}
\newcommand{\tnleanPrBase}{https://github.com/LionSR/TNLean/pull/}
\newcommand{\tnleanDocBase}{https://sirui-lu.com/TNLean/docs/}
\newcommand{\ghissue}[1]{\href{\tnleanIssueBase#1}{GitHub issue \##1}}
\newcommand{\ghpr}[1]{\href{\tnleanPrBase#1}{PR \##1}}
\newcommand{\doclink}[1]{\href{\tnleanDocBase#1.html}{\path{#1}}}

% Dirac notation and the identity operator, as in blueprint/src/macros/common.tex.
% \providecommand keeps note-local definitions authoritative.
\usepackage{dsfont}
\providecommand{\ket}[1]{|#1\rangle}
\providecommand{\bra}[1]{\langle #1|}
\providecommand{\braket}[2]{\langle #1 | #2 \rangle}
\providecommand{\ketbra}[2]{|#1\rangle\!\langle #2|}
\providecommand{\Id}{\mathds{1}}
\providecommand{\one}{\mathds{1}}
\providecommand{\C}{\mathbb{C}}
\providecommand{\MN}[1]{M_{#1}(\C)}
\providecommand{\MD}{M_D(\C)}

% External referencing for paper sources.  The build helper writes sanitized
% auxiliary files under build/paper-aux/ and excludes duplicated source labels.
\usepackage{xr}

% Import a generated paper auxiliary file with a caller-supplied label prefix.
% The candidate paths support builds from docs/paper-gaps/, the repository root,
% and build/paper-aux/ itself.
\newcommand{\importpaperaux}[2]{%
  \IfFileExists{../../build/paper-aux/#2.aux}{%
    \externaldocument[#1]{../../build/paper-aux/#2}%
  }{%
    \IfFileExists{build/paper-aux/#2.aux}{%
      \externaldocument[#1]{build/paper-aux/#2}%
    }{%
      \IfFileExists{#2.aux}{%
        \externaldocument[#1]{#2}%
      }{}%
    }%
  }%
}

% General external reference macros
\newcommand{\paperref}[2]{\ref{#1-#2}}
\newcommand{\papereqref}[2]{(\ref{#1-#2})}

% CPSV16 source (1606.00608 / MPDO-22-12-17-2.tex) external document and shortcuts
\importpaperaux{cpsv16-}{1606.00608/MPDO-22-12-17-2-tnlean}
\newcommand{\cpsvref}[1]{\paperref{cpsv16}{#1}}
\newcommand{\cpsveqref}[1]{\papereqref{cpsv16}{#1}}

% Verdict marker: \gapnote{<kind>}{<status>}. Kind is the policy
% classification (clarification, local-correction, scope-restriction,
% unfaithful, false-source, open-gap); status is open, wip (elimination
% underway), resolved, or historical. Machine-read by the site index;
% prints a verdict line.
\newcommand{\gapnote}[2]{\par\noindent\textbf{Verdict:} #1 (#2).\par}


\title{Power Sums in the BNT Weight Comparison}
\author{}
\date{2026-05-08}

\begin{document}
\maketitle

\section{The source lemma}

This note concerns the scalar power-sum step in the non-periodic Fundamental
Theorem proof for matrix product states in canonical form.  The source is
\cite[Lemma~\texttt{Lem:app\_simple}, lines~1155--1163]{Cirac2016MPDO_arXiv}.
It considers two finite families
\[
  (\lambda_{a,k})_{k=1}^{x_a},
  \qquad
  (\lambda_{b,k})_{k=1}^{x_b},
\]
ordered by modulus and phase, and assumes
\[
  \sum_{k=1}^{x_a}\lambda_{a,k}^{N}
  =
  \sum_{k=1}^{x_b}\lambda_{b,k}^{N},
  \qquad N\le \max\{x_a,x_b\}.        \tag{\texttt{eq:app\_aux}}
\]
The conclusion is \(x_a=x_b\) and equality of the ordered lists.  In the
equal-case corollary of the Fundamental Theorem, this lemma is applied after
the BNT representatives have been matched, to the coefficient families
\[
  (\mu_{j,q})_{q=1}^{r_{a,j}},
  \qquad
  (\nu_{j,q}e^{i\phi_j})_{q=1}^{r_{b,j}}.
\]

The formal statements do not use the source's ordering convention.  They state
the conclusion as equality of multisets.  They do, however, keep the relevant
nonzero-coefficient hypothesis: positive powers alone cannot count zero
entries.

\section{The source-facing formal route}

For a fixed BNT representative \(j\), the formal coefficient comparison produces
eventual equality of the power sums
\[
  \sum_{q=1}^{r_j^S}(\mu^S_{j,q})^N
  =
  \sum_{q=1}^{r_j^T}(\mu^T_{j,q})^N
\]
from BNT linear independence.  A finite signed sum of nonzero geometric
sequences which vanishes eventually vanishes at every exponent, so the equality
holds for all positive \(N\).  The formal proof then applies the
unequal-cardinality finite-range power-sum theorem.\footnote{\path{Matrix.sum_pow_eq_implies_card_eq_and_multiset_eq_of_le_max_card}.}
The endpoint used by the BNT comparison is the corresponding sector-weight
comparison lemma.\footnote{\path{MPSTensor.SectorWeightData.copies_eq_and_weight_multiset_eq_of_eventually_coeff_eq}.}

The conclusion is
\[
  r_j^S=r_j^T,
  \qquad
  \{\!\{\mu^S_{j,q}:1\le q\le r_j^S\}\!\}
  =
  \{\!\{\mu^T_{j,q}:1\le q\le r_j^T\}\!\}.
\]
This is the multiset form of the source's ordered-list conclusion.

\section{The alternative formal route}

An earlier formal proof reached the same endpoint by using the recurrence
behind finite sums of geometric sequences before applying Newton--Girard.  The
mathematical input is the following elementary lemma.  Let
\[
  F(N)=\sum_{\ell=1}^{m} c_\ell w_\ell^N,
  \qquad w_\ell\neq 0.
\]
If \(F(N)=0\) for all \(N\ge M\), then \(F(k)=0\) for every \(k\ge 0\).  Indeed,
subtracting \(w_1F(N)\) from \(F(N+1)\) gives
\[
  0
  =
  F(N+1)-w_1F(N)
  =
  \sum_{\ell=2}^{m} c_\ell(w_\ell-w_1)w_\ell^N
\]
for all \(N\ge M\).  This is a signed sum with one fewer term.  Induction on
\(m\) shows that the displayed sum is zero for every \(N\ge 0\), hence
\[
  F(N+1)=w_1F(N)\qquad\text{for all }N\ge 0.
\]
Thus \(F(k)=w_1^kF(0)\).  Since \(F(M)=0\) and \(w_1^M\neq 0\), one has
\(F(0)=0\), and therefore \(F(k)=0\) for all \(k\ge 0\).

For the two sector-weight families, this lemma is applied to the signed union
\[
  F(N)
  =
  \sum_{q=1}^{r_j^S}(\mu^S_{j,q})^N
  -
  \sum_{q=1}^{r_j^T}(\mu^T_{j,q})^N .
\]
The nonzero-weight hypothesis is exactly what permits the step \(w_1^M\neq 0\).
Since all weights in \leanid{SectorWeightData} are nonzero, the exponent-zero
power sum is
\[
  \sum_{q=1}^{r_j}\mu_{j,q}^{0}=r_j.
\]
Thus \(F(0)=0\) gives \(r_j^S=r_j^T\).  After transporting the right-hand index
set along this equality, the same extrapolation gives equality of the positive
power sums over a common index type:
\[
  \sum_{q=1}^{r_j^S}(\mu^S_{j,q})^N
  =
  \sum_{q=1}^{r_j^S}(\mu^T_{j,q})^N,
  \qquad N\ge 1.
\]
The proof then applies the same-cardinality Newton--Girard theorem to recover
the weight multiset.

This route was previously present as checked Lean code through the declarations
\[
\begin{gathered}
  \text{\leanid{MPSTensor.SectorWeightData.copies_eq_of_eventually_coeff_eq}},\\
  \text{\leanid{MPSTensor.SectorWeightData.eventually_coeff_eq_implies_all_pos_eq_after_copies}},\\
  \text{\leanid{MPSTensor.SectorWeightData.weight_multiset_eq_of_copies_eq_of_coeff_eq}}.
\end{gathered}
\]
Those declarations have been removed from the current source tree because the
proportional Fundamental Theorem route no longer uses this same-cardinality
detour.  The calculation above records the removed argument for comparison
with the paper, not a present Lean proof.

\section{Conclusion}

There is no mathematical gap in the removed alternative route under the
nonzero-weight hypotheses used in the former sector-weight data: the
extrapolation theorem proves the \(N=0\) identity, so the copy count is
justified.  Nevertheless, this is no longer a present Lean route.  The
source-facing theorem surface should use the unequal-cardinality power-sum
lemma, because that is the lemma invoked in the paper after BNT representative
matching.

There would be a genuine gap if one tried to count zero weights from positive
power sums alone.  For example, adjoining a zero to a finite family does not
change any positive power sum.  Hence a statement without a nonzero-coefficient
hypothesis can only recover the multiset after deleting zero entries, unless an
additional \(N=0\) identity is assumed.

\bibliographystyle{alpha}
\bibliography{references}

\end{document}
